\documentclass[12pt]{article}
\usepackage{mathrsfs}
\usepackage{graphicx}
\usepackage{amssymb,amsmath,amsthm}
\usepackage{epstopdf}
\usepackage{bm}
%\usepackage{feynmf}
\usepackage[pdftex,colorlinks,urlcolor=blue]{hyperref}
\pagestyle{plain}
%\usepackage{mcite}
\DeclareGraphicsRule{.tif}{png}{.png}{`convert #1 `basename #1
.tif`.png} \textwidth = 6.5 in \textheight = 9 in \oddsidemargin =
0.0 in \evensidemargin = 0.0 in \topmargin = 0.0 in \headheight =
0.0 in \headsep = 0.0 in
\parskip = 0.1in
\parindent = 0.1in
\linespread{1.0}
\def\ba{\begin{eqnarray}}
\def\ea{\end{eqnarray}}
\def\nn{\nonumber}
\def\L{\mathcal{L}}
\def\H{\mathcal{H}}
\def\M{\mathcal{M}}
\def\N{\mathcal{N}}
\def\P{\mathcal{P}}
\def\R{\mathcal{R}}
\def\O{\mathcal{O}}
\def\q{\mathbf{q}}
\def\D{\mathcal{D}}
\def\V{\mathcal{V}}
\def\:{\boldsymbol{:}}
\def\B{\boldsymbol{B}}
\def\Bbeta{\boldsymbol{\beta}}
\def\x{\mathbf{x}}
\def\u{\mathbf{u}}
\def\y{\mathbf{y}}
\def\p{\mathbf{p}}
\def\v{\mathbf{v}}
\def\n{\hat{\mathbf{n}}}
\def\dk#1{\stackrel{\circ}{#1}}
\def\ddk#1{\stackrel{\circ \circ}{#1}}
\def\dkt#1{\stackrel{\bullet}{#1}}
\def\k{\mathbf{k}}
\def\F{\mathcal{F}}
\def\Tr{\textrm{Tr}}
\def\Vol{\textrm{Vol}}
\def\slash{\makebox[0cm][l]{/}}
\long\def\symbolfootnote[#1]#2{\begingroup%
\def\thefootnote{\fnsymbol{footnote}}\footnote[#1]{#2}\endgroup}


\begin{document}




\begin{flushright}
Chris Gauthier
\today
\end{flushright}
\begin{center} 
\vglue .06in
{\Large \bf {$U(1)$ Gauge Field Around a Black Hole}}\\[.5in]

\end{center}
%+++++++++++++++++++++++++++++++++++++++++++++++++++++++++++++++++++++++++++++++
%+++++++++++++++++++++++++++++++++++++++++++++++++++++++++++++++++++++++++++++++
%+++++++++++++++++++++++++++++++++++++++++++++++++++++++++++++++++++++++++++++++
\section{Maxwell's Equations in Curved Space: Tree Level}
%+++++++++++++++++++++++++++++++++++++++++++++++++++++++++++++++++++++++++++++++
%+++++++++++++++++++++++++++++++++++++++++++++++++++++++++++++++++++++++++++++++
%+++++++++++++++++++++++++++++++++++++++++++++++++++++++++++++++++++++++++++++++
%+++++++++++++++++++++++++++++++++++++++++++++++++++++++++++++++++++++++++++++++
Recall that Maxwell's equations in curved space are written
%+++++++++++++++++++++++++++++++++++++++++++++++++++++++++++++++++++++++++++++++
\begin{gather}
\nabla_{\mu}
F^{\mu\nu}
=
\partial_{\mu}
F^{\mu\nu}
+
\Gamma^{\rho}_{\,\,\, \mu \rho}
F^{\mu \nu}
+
\Gamma^{\nu}_{\,\,\, \rho \mu}
F^{\mu \rho}
=
4 \pi j^{\nu}
\end{gather}
%+++++++++++++++++++++++++++++++++++++++++++++++++++++++++++++++++++++++++++++++
where as always $F^{\mu\nu} = \nabla^{\mu} A^{\nu} - \nabla^{\nu} A^{\mu} =  \partial^{\mu} A^{\nu} - \partial^{\nu} A^{\mu}$. In terms of the electric and magnetic field strengths. 
%+++++++++++++++++++++++++++++++++++++++++++++++++++++++++++++++++++++++++++++++
\begin{gather}
F_{\mu\nu}
=
\left[
\begin{array}{cccc}
0 & - E_{1} & - E_{2} & - E_{3}
\\
E_{1} & 0 & B_{3} & - B_{2}
\\
E_{2} & - B_{3} & 0 & B_{1}
\\
E_{3} & B_{2} & - B_{1} & 0
\end{array}
\right]
\end{gather}
%+++++++++++++++++++++++++++++++++++++++++++++++++++++++++++++++++++++++++++++++
So in terms of the potential, the vacuum maxwell equation become
%+++++++++++++++++++++++++++++++++++++++++++++++++++++++++++++++++++++++++++++++
\begin{gather}
\nabla_{\mu}
F^{\mu\nu}
=
\square A^{\nu}
-
\nabla_{\mu}
\nabla^{\nu}
A^{\mu}
\end{gather}
%+++++++++++++++++++++++++++++++++++++++++++++++++++++++++++++++++++++++++++++++













%+++++++++++++++++++++++++++++++++++++++++++++++++++++++++++++++++++++++++++++++
%+++++++++++++++++++++++++++++++++++++++++++++++++++++++++++++++++++++++++++++++
%+++++++++++++++++++++++++++++++++++++++++++++++++++++++++++++++++++++++++++++++
\subsection{Covariant Form}
%+++++++++++++++++++++++++++++++++++++++++++++++++++++++++++++++++++++++++++++++
%+++++++++++++++++++++++++++++++++++++++++++++++++++++++++++++++++++++++++++++++
%+++++++++++++++++++++++++++++++++++++++++++++++++++++++++++++++++++++++++++++++
In this paper we will write the maxwell equations for the covariant vector potential $A_{\mu}$. To get the maxwell equations into a more manageable form let's look at the equation piece by piece starting with the $\square A_{\nu}$ term. 
%+++++++++++++++++++++++++++++++++++++++++++++++++++++++++++++++++++++++++++++++
\begin{gather}
\square
A_{\nu}
=
\nabla_{\mu}
\nabla^{\mu}
A_{\nu}
=
g^{\mu \rho}
\nabla_{\mu}
\nabla_{\rho}
A_{\nu}
=
g^{\mu \rho}
\nabla_{\mu}
\left(
\partial_{\rho}
A_{\nu}
-
\Gamma^{\sigma}_{\rho\nu}
A_{\sigma}
\right)
\nn \\
=
g^{\mu \rho}
\left[
\partial_{\mu}
\partial_{\rho}
A_{\nu}
-
\partial_{\mu}
\left(
\Gamma^{\sigma}_{\rho\nu}
A_{\sigma}
\right)
-
\Gamma^{\lambda}_{\nu \mu}
\left(
\partial_{\rho}
A_{\lambda}
-
\Gamma^{\sigma}_{\rho\lambda}
A_{\sigma}
\right)
-
\Gamma^{\lambda}_{\rho \mu}
\left(
\partial_{\lambda}
A_{\nu}
-
\Gamma^{\sigma}_{\lambda\nu}
A_{\sigma}
\right)
\right]
\nn \\
=
g^{\mu \rho}
\left[
\partial_{\mu}
\partial_{\rho}
A_{\nu}
-
\partial_{\mu}
\Gamma^{\sigma}_{\rho\nu}
A_{\sigma}
-
\Gamma^{\sigma}_{\rho\nu}
\partial_{\mu}
A_{\sigma}
-
\Gamma^{\lambda}_{\nu \mu}
\partial_{\rho}
A_{\lambda}
+
\Gamma^{\lambda}_{\nu \mu}
\Gamma^{\sigma}_{\rho\lambda}
A_{\sigma}
-
\Gamma^{\lambda}_{\rho \mu}
\partial_{\lambda}
A_{\nu}
+
\Gamma^{\lambda}_{\rho \mu}
\Gamma^{\sigma}_{\lambda\nu}
A_{\sigma}
\right]
\nn \\
=
g^{\mu \rho}
\left[
\partial_{\mu}
\partial_{\rho}
A_{\nu}
-
2 \Gamma^{\sigma}_{\rho\nu}
\partial_{\mu}
A_{\sigma}
-
\Gamma^{\lambda}_{\rho \mu}
\partial_{\lambda}
A_{\nu}
+
\left(
-
\partial_{\mu}
\Gamma^{\sigma}_{\rho\nu}
+
\Gamma^{\lambda}_{\nu \mu}
\Gamma^{\sigma}_{\rho\lambda}
+
\Gamma^{\lambda}_{\rho \mu}
\Gamma^{\sigma}_{\lambda\nu}
\right)
A_{\sigma}
\right]
\nn \\
=
g^{\mu \rho}
\left[
\partial_{\mu}
\partial_{\rho}
A_{\nu}
-
2 \Gamma^{\sigma}_{\rho\nu}
\partial_{\mu}
A_{\sigma}
-
\Gamma^{\lambda}_{\rho \mu}
\partial_{\lambda}
A_{\nu}
\right]
-
g^{\mu \rho}
\partial_{\nu}
\Gamma^{\sigma}_{\rho\mu}
A_{\sigma}
\nn \\
-
g^{\mu \rho}
\left(
\partial_{\mu}
\Gamma^{\sigma}_{\rho\nu}
-
\partial_{\nu}
\Gamma^{\sigma}_{\rho\mu}
-
\Gamma^{\lambda}_{\nu \mu}
\Gamma^{\sigma}_{\rho\lambda}
-
\Gamma^{\lambda}_{\rho \mu}
\Gamma^{\sigma}_{\lambda\nu}
\right)
A_{\sigma}
=
g^{\mu \rho}
\left[
\partial_{\mu}
\partial_{\rho}
A_{\nu}
-
2 \Gamma^{\sigma}_{\rho\nu}
\partial_{\mu}
A_{\sigma}
-
\Gamma^{\lambda}_{\rho \mu}
\partial_{\lambda}
A_{\nu}
\right]
\nn \\
-
\partial_{\nu}
(g^{\mu \rho}
\Gamma^{\sigma}_{\rho\mu})
A_{\sigma}
+
\Gamma^{\sigma}_{\rho\mu}
\partial_{\nu}
g^{\mu \rho}
A_{\sigma}
-
g^{\mu \rho}
\left(
\partial_{\mu}
\Gamma^{\sigma}_{\rho\nu}
-
\partial_{\nu}
\Gamma^{\sigma}_{\rho\mu}
-
\Gamma^{\lambda}_{\nu \mu}
\Gamma^{\sigma}_{\rho\lambda}
-
\Gamma^{\lambda}_{\rho \mu}
\Gamma^{\sigma}_{\lambda\nu}
\right)
A_{\sigma}
\end{gather}
%+++++++++++++++++++++++++++++++++++++++++++++++++++++++++++++++++++++++++++++++
A short detour
%+++++++++++++++++++++++++++++++++++++++++++++++++++++++++++++++++++++++++++++++
\begin{gather}
0
=
\partial_{\nu}\eta^{\mu}_{\,\,\, \sigma}
=
\partial_{\nu}(g^{\mu \lambda}g_{\lambda\sigma})
=
g_{\lambda \sigma}
\partial_{\nu}
g^{\mu \lambda}
+
g^{\mu \lambda}
\partial_{\nu}
g_{\lambda \sigma}
\nn \\
\Rightarrow
\quad
0
=
g^{\sigma \rho}
g_{\lambda \sigma}
\partial_{\nu}
g^{\mu \lambda}
+
g^{\sigma \rho}
g^{\mu \lambda}
\partial_{\nu}
g_{\lambda \sigma}
=
\eta^{\rho}_{\,\,\, \lambda}
\partial_{\nu}
g^{\mu \lambda}
+
g^{\sigma \rho}
g^{\mu \lambda}
\partial_{\nu}
g_{\lambda \sigma}
\nn \\
=
\partial_{\nu}
g^{\mu \rho}
+
g^{\sigma \rho}
g^{\mu \lambda}
\partial_{\nu}
g_{\lambda \sigma}
\quad
\Rightarrow
\quad
\partial_{\nu}
g^{\mu \rho}
=-
g^{\omega \rho}
g^{\mu \lambda}
\partial_{\nu}
g_{\lambda \omega}
\end{gather}
%+++++++++++++++++++++++++++++++++++++++++++++++++++++++++++++++++++++++++++++++
continuing
%+++++++++++++++++++++++++++++++++++++++++++++++++++++++++++++++++++++++++++++++
\begin{gather}
\square
A_{\nu}
=
g^{\mu \rho}
\left[
\partial_{\mu}
\partial_{\rho}
A_{\nu}
-
2 \Gamma^{\sigma}_{\rho\nu}
\partial_{\mu}
A_{\sigma}
-
\Gamma^{\lambda}_{\rho \mu}
\partial_{\lambda}
A_{\nu}
\right]
\nn \\
-
\partial_{\nu}
(g^{\mu \rho}
\Gamma^{\sigma}_{\rho\mu})
A_{\sigma}
-
\Gamma^{\sigma}_{\rho\mu}
g^{\omega \rho}
g^{\mu \lambda}
\partial_{\nu}
g_{\lambda \omega}
A_{\sigma}
-
g^{\mu \rho}
\left(
\partial_{\mu}
\Gamma^{\sigma}_{\rho\nu}
-
\partial_{\nu}
\Gamma^{\sigma}_{\rho\mu}
-
\Gamma^{\lambda}_{\nu \mu}
\Gamma^{\sigma}_{\rho\lambda}
-
\Gamma^{\lambda}_{\rho \mu}
\Gamma^{\sigma}_{\lambda\nu}
\right)
A_{\sigma}
\nn \\
=
g^{\mu \rho}
\left[
\partial_{\mu}
\partial_{\rho}
A_{\nu}
-
2 \Gamma^{\sigma}_{\rho\nu}
\partial_{\mu}
A_{\sigma}
-
\Gamma^{\lambda}_{\rho \mu}
\partial_{\lambda}
A_{\nu}
\right]
\nn \\
-
\partial_{\nu}
(g^{\mu \rho}
\Gamma^{\sigma}_{\rho\mu})
A_{\sigma}
-
\Gamma^{\sigma}_{\lambda \rho}
g^{\omega \lambda}
g^{\rho\mu}
\left(
\partial_{\nu}
g_{\mu \omega}
+
\partial_{\mu}
g_{\nu \omega}
-
\partial_{\omega}
g_{\mu \nu}
\right)
A_{\sigma}
\nn \\
-
g^{\mu \rho}
\left(
\partial_{\mu}
\Gamma^{\sigma}_{\rho\nu}
-
\partial_{\nu}
\Gamma^{\sigma}_{\rho\mu}
-
\Gamma^{\lambda}_{\nu \mu}
\Gamma^{\sigma}_{\rho\lambda}
-
\Gamma^{\lambda}_{\rho \mu}
\Gamma^{\sigma}_{\lambda\nu}
\right)
A_{\sigma}
\nn \\
=
g^{\mu \rho}
\left[
\partial_{\mu}
\partial_{\rho}
A_{\nu}
-
2 \Gamma^{\sigma}_{\rho\nu}
\partial_{\mu}
A_{\sigma}
-
\Gamma^{\lambda}_{\rho \mu}
\partial_{\lambda}
A_{\nu}
\right]
\nn \\
-
\partial_{\nu}
(g^{\mu \rho}
\Gamma^{\sigma}_{\rho\mu})
A_{\sigma}
-
2 \Gamma^{\sigma}_{\lambda \rho}
\Gamma^{\lambda}_{\mu\nu}
g^{\rho\mu}
A_{\sigma}
-
g^{\mu \rho}
\left(
\partial_{\mu}
\Gamma^{\sigma}_{\rho\nu}
-
\partial_{\nu}
\Gamma^{\sigma}_{\rho\mu}
-
\Gamma^{\lambda}_{\nu \mu}
\Gamma^{\sigma}_{\rho\lambda}
-
\Gamma^{\lambda}_{\rho \mu}
\Gamma^{\sigma}_{\lambda\nu}
\right)
A_{\sigma}
\nn \\
=
g^{\mu \rho}
\left[
\partial_{\mu}
\partial_{\rho}
A_{\nu}
-
2 \Gamma^{\sigma}_{\rho\nu}
\partial_{\mu}
A_{\sigma}
-
\Gamma^{\lambda}_{\rho \mu}
\partial_{\lambda}
A_{\nu}
\right]
\nn \\
-
\partial_{\nu}
(g^{\mu \rho}
\Gamma^{\sigma}_{\rho\mu})
A_{\sigma}
-
g^{\mu \rho}
\left(
\partial_{\mu}
\Gamma^{\sigma}_{\nu\rho}
-
\partial_{\nu}
\Gamma^{\sigma}_{\mu\rho}
+
\Gamma^{\sigma}_{\rho\lambda}
\Gamma^{\lambda}_{\nu \mu}
-
\Gamma^{\sigma}_{\nu \lambda}
\Gamma^{\lambda}_{\rho \mu}
\right)
A_{\sigma}
\nn \\
=
g^{\mu \rho}
\left[
\partial_{\mu}
\partial_{\rho}
A_{\nu}
-
2 \Gamma^{\sigma}_{\rho\nu}
\partial_{\mu}
A_{\sigma}
-
\Gamma^{\lambda}_{\rho \mu}
\partial_{\lambda}
A_{\nu}
\right]
\nn \\
-
\partial_{\nu}
(g^{\mu \rho}
\Gamma^{\sigma}_{\rho\mu})
A_{\sigma}
-
g^{\mu \rho}
\left(
\partial_{\mu}
\Gamma^{\sigma}_{\nu\rho}
-
\partial_{\nu}
\Gamma^{\sigma}_{\mu\rho}
+
\Gamma^{\sigma}_{\mu \lambda}
\Gamma^{\lambda}_{\nu \rho}
-
\Gamma^{\sigma}_{\nu \lambda}
\Gamma^{\lambda}_{\mu \rho}
\right)
A_{\sigma}
\nn \\
=
g^{\mu \rho}
\left[
\partial_{\mu}
\partial_{\rho}
A_{\nu}
-
2 \Gamma^{\sigma}_{\rho\nu}
\partial_{\mu}
A_{\sigma}
-
\Gamma^{\lambda}_{\rho \mu}
\partial_{\lambda}
A_{\nu}
\right]
-
\partial_{\nu}
(g^{\mu \rho}
\Gamma^{\sigma}_{\rho\mu})
A_{\sigma}
-
g^{\mu \rho}
R^{\sigma}_{\,\,\, \rho \mu \nu}
A_{\sigma}
\end{gather}
%+++++++++++++++++++++++++++++++++++++++++++++++++++++++++++++++++++++++++++++++
Using the symmetry properties of the curvature tensor
%+++++++++++++++++++++++++++++++++++++++++++++++++++++++++++++++++++++++++++++++
\begin{gather}
g^{\mu \rho}
R^{\sigma}_{\,\,\, \rho \mu \nu}
=
g^{\sigma \lambda}
g^{\mu \rho}
R_{\lambda \rho \mu \nu}
=
-
g^{\sigma \lambda}
g^{\mu \rho}
R_{\rho \lambda \mu \nu}
=
-
g^{\sigma \lambda}
R_{\lambda \nu}
\end{gather}
%+++++++++++++++++++++++++++++++++++++++++++++++++++++++++++++++++++++++++++++++
Thus, $\square A_{\nu}$ becomes:
%+++++++++++++++++++++++++++++++++++++++++++++++++++++++++++++++++++++++++++++++
\begin{gather}
\square
A_{\nu}
=
L A_{\nu}
-
2 g^{\mu \rho} \Gamma^{\sigma}_{\rho\nu}
\partial_{\mu}
A_{\sigma}
-
\partial_{\nu}
(g^{\mu \rho}
\Gamma^{\sigma}_{\rho\mu})
A_{\sigma}
+
g^{\sigma \mu}
R_{\mu \nu}
A_{\sigma}
\end{gather}
%+++++++++++++++++++++++++++++++++++++++++++++++++++++++++++++++++++++++++++++++
where $L$ is the scalar D'Alembert operator defined by
%+++++++++++++++++++++++++++++++++++++++++++++++++++++++++++++++++++++++++++++++
\begin{gather}
L
A_{\nu}
=
\frac{1}{\sqrt{- g}}
\partial_{\mu}(\sqrt{-g} g^{\mu\rho} \partial_{\rho} A_{\nu})
=
g^{\mu \rho}
\left[
\partial_{\mu}
\partial_{\rho}
A_{\nu}
-
\Gamma^{\lambda}_{\rho \mu}
\partial_{\lambda}
A_{\nu}
\right]
\end{gather}
%+++++++++++++++++++++++++++++++++++++++++++++++++++++++++++++++++++++++++++++++
Thus, Maxwell's equation becomes
%+++++++++++++++++++++++++++++++++++++++++++++++++++++++++++++++++++++++++++++++
\begin{gather}
\square A_{\nu}
-
\nabla_{\mu}
\nabla_{\nu}
A^{\mu}
=
\square A_{\nu}
-
\nabla_{\nu}
\nabla_{\mu}
A^{\mu}
-
R_{\nu \mu}
A^{\mu}
\nn \\
=
L A_{\nu}
-
2 g^{\mu \rho} \Gamma^{\sigma}_{\rho\nu}
\partial_{\mu}
A_{\sigma}
-
\partial_{\nu}
(g^{\mu \rho}
\Gamma^{\sigma}_{\rho\mu})
A_{\sigma}
+
g^{\sigma \mu}
R_{\mu \nu}
A_{\sigma}
-
\nabla_{\nu}
\nabla_{\mu}
A^{\mu}
-
R_{\nu \mu }
A^{\mu}
\nn \\
=
L A_{\nu}
-
2 g^{\mu \rho} \Gamma^{\sigma}_{\rho\nu}
\partial_{\mu}
A_{\sigma}
-
\partial_{\nu}
(g^{\mu \rho}
\Gamma^{\sigma}_{\rho\mu})
A_{\sigma}
+
R_{\mu \nu}
A^{\mu}
-
\nabla_{\nu}
\nabla_{\mu}
A^{\mu}
-
R_{\mu \nu}
A^{\mu}
\nn \\
=
L A_{\nu}
-
2 g^{\mu \rho} \Gamma^{\sigma}_{\rho\nu}
\partial_{\mu}
A_{\sigma}
-
\partial_{\nu}
(g^{\mu \rho}
\Gamma^{\sigma}_{\rho\mu})
A_{\sigma}
-
\nabla_{\nu}
\nabla_{\mu}
A^{\mu}
\nn \\
\Rightarrow
\quad
L A_{\nu}
=
2 g^{\mu \rho} \Gamma^{\sigma}_{\rho\nu}
\partial_{\mu}
A_{\sigma}
+
\partial_{\nu}
(g^{\mu \rho}
\Gamma^{\sigma}_{\rho\mu})
A_{\sigma}
+
\partial_{\nu}
(\nabla_{\mu}
A^{\mu})
\end{gather}
%+++++++++++++++++++++++++++++++++++++++++++++++++++++++++++++++++++++++++++++++




%+++++++++++++++++++++++++++++++++++++++++++++++++++++++++++++++++++++++++++++++
%+++++++++++++++++++++++++++++++++++++++++++++++++++++++++++++++++++++++++++++++
%+++++++++++++++++++++++++++++++++++++++++++++++++++++++++++++++++++++++++++++++
\subsection{Schwarzschild Space with Stereographic Coordinates}
%+++++++++++++++++++++++++++++++++++++++++++++++++++++++++++++++++++++++++++++++
%+++++++++++++++++++++++++++++++++++++++++++++++++++++++++++++++++++++++++++++++
%+++++++++++++++++++++++++++++++++++++++++++++++++++++++++++++++++++++++++++++++
The schwarzschild metric is given by
%+++++++++++++++++++++++++++++++++++++++++++++++++++++++++++++++++++++++++++++++
\begin{gather}
ds^{2}
=
(1 - \frac{2 M}{r})
dt^{2}
-
(1 - \frac{2 M}{r})^{-1}
dr^{2}
-
r^{2}
d\theta^{2}
-
r^{2} \sin^{2}\theta
d\varphi^{2}
\end{gather}
%+++++++++++++++++++++++++++++++++++++++++++++++++++++++++++++++++++++++++++++++
It is more covinent to use stereographic coordinates $(z,\bar{z})$ in place of the angular coordinates $(\theta,\varphi)$. The definition of the stereographic coordinates $(z,\bar{z})$ in terms of $(\theta,\varphi)$ is given by
%+++++++++++++++++++++++++++++++++++++++++++++++++++++++++++++++++++++++++++++++
\begin{gather}
z
=
\frac{\sin \theta e^{i \varphi}}{1 - \cos \theta}
,
\quad
\bar{z}
=
\frac{\sin \theta e^{-i \varphi}}{1 - \cos \theta}
\end{gather}
%+++++++++++++++++++++++++++++++++++++++++++++++++++++++++++++++++++++++++++++++
The metric in terms of these stereographic coordinates is
%+++++++++++++++++++++++++++++++++++++++++++++++++++++++++++++++++++++++++++++++
\begin{gather}
ds^{2}
=
(1 - \frac{2 M}{r}) dt^{2}
-
(1 - \frac{2 M}{r})^{-1}
dr^{2}
-
\frac{4 r^{2}}{(1 + z \bar{z})^{2}}
dz
d\bar{z}
\end{gather}
%+++++++++++++++++++++++++++++++++++++++++++++++++++++++++++++++++++++++++++++++
The Christoffel terms are then
%+++++++++++++++++++++++++++++++++++++++++++++++++++++++++++++++++++++++++++++++
\begin{gather}
g^{\mu\rho} \Gamma^{\sigma}_{\rho 0}
=
\Gamma^{\sigma \mu}_{\,\,\,\,\,\,\, 0}
=
\left[
\begin{array}{cccc}
0 & -\frac{M}{r^{2}} & 0 & 0
\\
\frac{M}{r^{2}} & 0 & 0 & 0
\\
0 & 0 & 0 & 0
\\
0 & 0 & 0 & 0
\end{array}
\right]
\nn \\
\Gamma^{\sigma \mu}_{\,\,\,\,\,\,\, 1}
=
\left[
\begin{array}{cccc}
\frac{M}{r^{2} (1 - \frac{2 M}{r})^{2}} & 0 & 0 & 0
\\
0 & \frac{M}{r^{2}} & 0 & 0
\\
0 & 0 & 0 & - \frac{(1 + z \bar{z})^{2}}{2 r^{3}}
\\
0 & 0 & - \frac{(1 + z \bar{z})^{2}}{2 r^{3}} & 0
\end{array}
\right]
\nn \\
\Gamma^{\sigma \mu}_{\,\,\,\,\,\,\, 2}
=
\left[
\begin{array}{cccc}
0 & 0 & 0 & 0
\\
0 & 0 & \frac{1}{r} (1 - \frac{2 M}{r}) & 0
\\
0 & -\frac{1}{r} (1 - \frac{2 M}{r}) & 0 & \frac{\bar{z} (1 + z \bar{z})}{r^{2}}
\\
0 & 0 & 0 & 0
\end{array}
\right]
\nn \\
\Gamma^{\sigma \mu}_{\,\,\,\,\,\,\, 3}
=
\left[
\begin{array}{cccc}
0 & 0 & 0 & 0
\\
0 & 0 & 0 & \frac{1}{r} (1 - \frac{2 M}{r}) 
\\
0 & 0 & 0 & 0
\\
0 & - \frac{1}{r} (1 - \frac{2 M}{r}) &  \frac{z (1 + z \bar{z})}{r^{2}}& 0
\end{array}
\right]
\end{gather}
%+++++++++++++++++++++++++++++++++++++++++++++++++++++++++++++++++++++++++++++++
and
%+++++++++++++++++++++++++++++++++++++++++++++++++++++++++++++++++++++++++++++++
\begin{gather}
g^{\mu\nu}
\Gamma^{\rho}_{\mu\nu}
=
\left\{
0, \frac{2 (r - M)}{r^{2}} , 0 , 0
\right\}
\nn \\
\partial_{0}\left(g^{\mu\nu}
\Gamma^{\rho}_{\mu\nu}\right)
=
\left\{
0,0,0,0
\right\}
,
\quad
\partial_{1}\left(g^{\mu\nu}
\Gamma^{\rho}_{\mu\nu}\right)
=
\left\{
0,- \frac{2}{r^{2}}(1 - \frac{2 M}{r}),0,0
\right\}
,
\nn \\
\partial_{2}\left(g^{\mu\nu}
\Gamma^{\rho}_{\mu\nu}\right)
=
\left\{
0, 0 , 0 ,0
\right\}
,
\quad
\partial_{3}\left(g^{\mu\nu}
\Gamma^{\rho}_{\mu\nu}\right)
=
\left\{
0,0,0,0
\right\}
\end{gather}
%+++++++++++++++++++++++++++++++++++++++++++++++++++++++++++++++++++++++++++++++

%+++++++++++++++++++++++++++++++++++++++++++++++++++++++++++++++++++++++++++++++
\begin{gather}
\Gamma^{\sigma \mu}_{\,\,\,\,\,\,\, 0}
\partial_{\mu}
A_{\sigma}
=
\frac{M}{r^{2}}
\left(
\partial_{0}A_{1}
-
\partial_{1}A_{0}
\right)
,
\nn \\
\Gamma^{\sigma \mu}_{\,\,\,\,\,\,\, 1}
\partial_{\mu}
A_{\sigma}
=
\frac{M}{r^{2} (1 - \frac{2 M}{r})^{2}} 
\partial_{0}
A_{0}
+
\frac{M}{r^{2}}
\partial_{1}
A_{1}
-
\frac{(1 + z \bar{z})^{2}}{2 r^{3}}
(
\partial_{2}
A_{3}
+
\partial_{3}
A_{2})
,
\nn \\
\Gamma^{\sigma \mu}_{\,\,\,\,\,\,\, 2}
\partial_{\mu}
A_{\sigma}
=
-
\frac{1}{r}
(1 - \frac{2 M}{r})
(
\partial_{1}
A_{2}
-
\partial_{2}
A_{1}
)
+
\frac{\bar{z} (1 + z \bar{z})}{r^{2}}
\partial_{3}
A_{2}
,
\nn \\
\Gamma^{\sigma \mu}_{\,\,\,\,\,\,\, 3}
\partial_{\mu}
A_{\sigma}
=
-
\frac{1}{r}
(1 - \frac{2 M}{r})
(
\partial_{1}
A_{3}
-
\partial_{3}
A_{1}
)
+
\frac{z (1 + z \bar{z})}{r^{2}}
\partial_{2}
A_{3}
\end{gather}
%+++++++++++++++++++++++++++++++++++++++++++++++++++++++++++++++++++++++++++++++


%+++++++++++++++++++++++++++++++++++++++++++++++++++++++++++++++++++++++++++++++
\begin{gather}
\partial_{0}\left(g^{\mu\nu}
\Gamma^{\rho}_{\mu\nu}\right)
A_{\rho}
=
0
,
\quad
\partial_{1}\left(g^{\mu\nu}
\Gamma^{\rho}_{\mu\nu}\right)
=
- \frac{2}{r^{2}}(1 - \frac{2 M}{r}) A_{1},
\nn \\
\partial_{2}\left(g^{\mu\nu}
\Gamma^{\rho}_{\mu\nu}\right)
A_{\rho}
=
0
,
\quad
\partial_{3}\left(g^{\mu\nu}
\Gamma^{\rho}_{\mu\nu}\right)
A_{\rho}
=
0
\end{gather}
%+++++++++++++++++++++++++++++++++++++++++++++++++++++++++++++++++++++++++++++++
So the equations of motion
%+++++++++++++++++++++++++++++++++++++++++++++++++++++++++++++++++++++++++++++++
\begin{gather}
L A_{0}
=
\frac{2 M}{r^{2}}
\left(
\partial_{0}A_{1}
-
\partial_{1}A_{0}
\right)
+
\partial_{0}
\nabla_{\mu}
A^{\mu}
\label{MaxEqu1}
\\
L A_{1}
=
\frac{2 M}{r^{2} (1 - \frac{2 M}{r})^{2}} 
\partial_{0}
A_{0}
+
\frac{2 M}{r^{2}}
\partial_{1}
A_{1}
\nn \\
-
\frac{(1 + z \bar{z})^{2}}{r^{3}}
(
\partial_{2}
A_{3}
+
\partial_{3}
A_{2})
- 
\frac{2}{r^{2}}(1 - \frac{2 M}{r}) A_{1}
+
\partial_{1}
\nabla_{\mu}
A^{\mu}
\label{MaxEqu2}
\\
L A_{2}
=
-
\frac{2}{r}
(1 - \frac{2 M}{r})
(
\partial_{1}
A_{2}
-
\partial_{2}
A_{1}
)
+
\frac{2 \bar{z} (1 + z \bar{z})}{r^{2}}
\partial_{3}
A_{2}
+
\partial_{2}
\nabla_{\mu}
A^{\mu}
\label{MaxEqu3}
\\
L A_{3}
=
-
\frac{2}{r}
(1 - \frac{2 M}{r})
(
\partial_{1}
A_{3}
-
\partial_{3}
A_{1}
)
+
\frac{2 z (1 + z \bar{z})}{r^{2}}
\partial_{2}
A_{3}
+
\partial_{3}
\nabla_{\mu}
A^{\mu}
\label{MaxEqu4}
\end{gather}
%+++++++++++++++++++++++++++++++++++++++++++++++++++++++++++++++++++++++++++++++
Expand out $\nabla_{\mu} A^{\mu}$
%+++++++++++++++++++++++++++++++++++++++++++++++++++++++++++++++++++++++++++++++
\begin{gather}
\nabla_{\mu}A^{\mu}
=
g^{\mu\nu} \nabla_{\mu}A_{\nu}
=
g^{\mu\nu}
\left(
\partial_{\mu}
A_{\nu}
-
\Gamma^{\lambda}_{\mu \nu}
A_{\lambda}
\right)
\nn \\
=
(1 - \frac{2 M}{r})^{-1}
\partial_{0}
A_{0}
-
(1 - \frac{2 M}{r})
\partial_{1}
A_{1}
-
\frac{(1 + z \bar{z})^{2}}{2 r^{2}}
(\partial_{2} A_{3} + \partial_{3} A_{2})
-
\frac{2 (r - M)}{r^{2}}
A_{1}
\nn \\
=
(1 - \frac{2 M}{r})^{-1}
\partial_{0}
A_{0}
-
(1 - \frac{2 M}{r})
\frac{1}{r^{2}}
\partial_{1}
(r^{2}
A_{1})
-
\frac{2 M}{r^{2}}
A_{1}
-
\frac{(1 + z \bar{z})^{2}}{2 r^{2}}
(\partial_{2} A_{3} + \partial_{3} A_{2})
\end{gather}
%+++++++++++++++++++++++++++++++++++++++++++++++++++++++++++++++++++++++++++++++
Fix the gauge so that $A_{0} =0$, then equation \ref{MaxEqu1} becomes
%+++++++++++++++++++++++++++++++++++++++++++++++++++++++++++++++++++++++++++++++
\begin{gather}
0
=
\partial_{0}
\left[
\frac{2 M}{r^{2}}
A_{1}
-
(1 - \frac{2 M}{r})
\partial_{1}
A_{1}
-
\frac{(1 + z \bar{z})^{2}}{2 r^{2}}
(\partial_{2} A_{3} + \partial_{3} A_{2})
-
\frac{2 (r - M)}{r^{2}}
A_{1}
\right]
\nn \\
=
-
\partial_{0}
\left[
(1 - \frac{2 M}{r})
\partial_{1}
A_{1}
+
\frac{(1 + z \bar{z})^{2}}{2 r^{2}}
(\partial_{2} A_{3} + \partial_{3} A_{2})
+
\frac{2}{r}
 (1 - \frac{2 M}{r})
A_{1}
\right]
\nn \\
=
-
\partial_{0}
\left[
(1 - \frac{2 M}{r})
\frac{1}{r^{2}}
\partial_{1}
(r^{2}
A_{1})
+
\frac{(1 + z \bar{z})^{2}}{2 r^{2}}
(\partial_{2} A_{3} + \partial_{3} A_{2})
\right]
\end{gather}
%+++++++++++++++++++++++++++++++++++++++++++++++++++++++++++++++++++++++++++++++
Thus $\nabla_{\mu} A^{\mu}$ becomes in this gauge
%+++++++++++++++++++++++++++++++++++++++++++++++++++++++++++++++++++++++++++++++
\begin{gather}
\nabla_{\mu}A^{\mu}
=
-
\frac{2 M}{r^{2}}
A_{1}
\end{gather}
%+++++++++++++++++++++++++++++++++++++++++++++++++++++++++++++++++++++++++++++++
Thus equation (\ref{MaxEqu2})
%+++++++++++++++++++++++++++++++++++++++++++++++++++++++++++++++++++++++++++++++
\begin{gather}
L A_{1}
=
\frac{1}{r}
\left[
-
(1-
\frac{2 M}{r})
\frac{1}{r^{2}}
\partial_{1}(r^{2} A_{1})
+
\partial_{1}
A_{1}
-
\frac{(1 + z \bar{z})^{2}}{r^{2}}
(
\partial_{2}
A_{3}
+
\partial_{3}
A_{2})
+
r
\partial_{1}
\nabla_{\mu}
A^{\mu}
\right]
\nn \\
\Rightarrow
\quad
L A_{1}
=
\frac{1}{r}
\left[
(1-
\frac{2 M}{r})
\frac{1}{r^{2}}
\partial_{1}(r^{2} A_{1})
+
\partial_{1}
A_{1}
-
2 M r
\partial_{1}
\left(
\frac{A_{1}}{r^{2}}
\right)
\right]
\nn \\
\Rightarrow
\quad
L A_{1}
=
\frac{1}{r}
\left[
(1-
\frac{2 M}{r})
\frac{1}{r^{2}}
\partial_{1}(r^{2} A_{1})
+
\partial_{1}
A_{1}
+
\frac{4 M}{r^{2}}
A_{1}
-
\frac{2 M}{r}
\partial_{1}
A_{1}
\right]
\nn \\
\Rightarrow
\quad
L A_{1}
=
\frac{1}{r}
\left[
2
(1 - \frac{2 M}{r})
\partial_{1}
A_{1}
+
\frac{2}{r}
(1-
\frac{2 M}{r})
A_{1}
+
\frac{4 M}{r^{2}}
A_{1}
\right]
\nn \\
\Rightarrow
\quad
L A_{1}
=
\frac{2}{r}
\left[
(1 - \frac{2 M}{r})
\partial_{1}
A_{1}
+
\frac{1}{r}
(1-
\frac{2 M}{r})
A_{1}
+
\frac{2 M}{r^{2}}
A_{1}
\right]
\nn \\
\Rightarrow
\quad
L A_{1}
=
\frac{2}{r}
\left[
(1 - \frac{2 M}{r})
\partial_{1}
A_{1}
+
\frac{1}{r}
\partial_{1}
\left(
r
(1-
\frac{2 M}{r})
\right)
A_{1}
\right]
\nn \\
\Rightarrow
\quad
L A_{r}
=
\frac{2}{r^{2}}
\partial_{r}
\left[
r
(1-
\frac{2 M}{r})
A_{r}
\right]
\end{gather}
%+++++++++++++++++++++++++++++++++++++++++++++++++++++++++++++++++++++++++++++++
where $A_{r} = A_{1}$ and $\partial_{r} = \partial_{1}$. Equation (\ref{MaxEqu3})
%+++++++++++++++++++++++++++++++++++++++++++++++++++++++++++++++++++++++++++++++
\begin{gather}
L A_{2}
=
-
\frac{2}{r}
(1 - \frac{2 M}{r})
(
\partial_{1}
A_{2}
-
\partial_{2}
A_{1}
)
+
\frac{2 \bar{z} (1 + z \bar{z})}{r^{2}}
\partial_{3}
A_{2}
-
\frac{2 M}{r^{2}}
\partial_{2}
A_{1}
\nn \\
\Rightarrow
\quad
\left(L 
-
\frac{2 \bar{z} (1 + z \bar{z})}{r^{2}}
\partial_{\bar{z}}
\right)
A_{z}
=
-
\frac{2}{r}
(1 - \frac{2 M}{r})
\partial_{r}
A_{z}
+
\frac{2}{r}
(1 - \frac{3 M}{r})
\partial_{z}
A_{r}
\label{MaxEqu3Two}
\end{gather}
%+++++++++++++++++++++++++++++++++++++++++++++++++++++++++++++++++++++++++++++++
where, $A_{z} = A_{2}$, $\partial_{z} = \partial_{2}$ and $\partial_{\bar{z}} = \partial_{3}$. Equation (\ref{MaxEqu4}) becomes
%+++++++++++++++++++++++++++++++++++++++++++++++++++++++++++++++++++++++++++++++
\begin{gather}
L A_{3}
=
-
\frac{2}{r}
(1 - \frac{2 M}{r})
(
\partial_{1}
A_{3}
-
\partial_{3}
A_{1}
)
+
\frac{2 z (1 + z \bar{z})}{r^{2}}
\partial_{2}
A_{3}
-
\frac{2 M}{r^{2}}
\partial_{3}
A_{1}
\nn \\
\Rightarrow
\quad
\left(
L 
-
\frac{2 z (1 + z \bar{z})}{r^{2}}
\partial_{z}
\right)
A_{\bar{z}}
=
-
\frac{2}{r}
(1 - \frac{2 M}{r})
\partial_{r}
A_{\bar{z}}
+
\frac{2}{r}
(1 - \frac{3 M}{r})
\partial_{\bar{z}}
A_{r}
\end{gather}
%+++++++++++++++++++++++++++++++++++++++++++++++++++++++++++++++++++++++++++++++
where $A_{\bar{z}} = A_{3}$. Let's concentrate on the left hand side of (\ref{MaxEqu3Two})
%+++++++++++++++++++++++++++++++++++++++++++++++++++++++++++++++++++++++++++++++
\begin{gather}
\left(
L
-
\frac{2 \bar{z} (1 + z \bar{z})}{r^{2}}
\partial_{\bar{z}}
\right)
\partial_{z}
\psi
\end{gather}
%+++++++++++++++++++++++++++++++++++++++++++++++++++++++++++++++++++++++++++++++
Looking in particular at the $L \partial_{z}\psi$ piece 
%+++++++++++++++++++++++++++++++++++++++++++++++++++++++++++++++++++++++++++++++
\begin{gather}
L\partial_{z}
\psi
=
\frac{1}{\sqrt{- g}}
\partial_{\mu}
(\sqrt{-g} g^{\mu\nu}\partial_{\nu}\partial_{z} \psi)
=
\frac{1}{\sqrt{- g}}
\partial_{\mu}
\left[
\partial_{z}(\sqrt{-g} g^{\mu\nu}\partial_{\nu} \psi)
-
\partial_{z}(\sqrt{-g} g^{\mu\nu}) \partial_{\nu} \psi
\right]
\nn \\
=
\frac{1}{\sqrt{- g}}
\partial_{\mu} \partial_{z}(\sqrt{-g} g^{\mu\nu}\partial_{\nu} \psi)
-
\frac{1}{\sqrt{- g}}
\partial_{\mu} \left[
\partial_{z}(\sqrt{-g} g^{\mu\nu}) \partial_{\nu} \psi\right]
\nn \\
=
\partial_{z}
\left[
\frac{1}{\sqrt{- g}}
\partial_{\mu} (\sqrt{-g} g^{\mu\nu}\partial_{\nu} \psi)
\right]
-
\partial_{z}
\left[
\frac{1}{\sqrt{- g}}
\right]
\partial_{\mu} (\sqrt{-g} g^{\mu\nu}\partial_{\nu} \psi)
-
\frac{1}{\sqrt{- g}}
\partial_{\mu} \left[
\partial_{z}(\sqrt{-g} g^{\mu\nu}) \partial_{\nu} \psi\right]
\nn \\
=
\partial_{z}
L\psi
+
\frac{1}{-g}
\partial_{z}
\left[
\sqrt{- g}
\right]
\partial_{\mu} (\sqrt{-g} g^{\mu\nu}\partial_{\nu} \psi)
-
\frac{1}{\sqrt{- g}}
\partial_{\mu} \left[
\partial_{z}(\sqrt{-g} g^{\mu\nu}) \partial_{\nu} \psi\right]
\nn \\
=
\partial_{z}
L\psi
+
\frac{1}{-g}
\partial_{z}
\left[
\sqrt{- g}
\right]
\partial_{\mu} (\sqrt{-g} g^{\mu\nu}) \partial_{\nu} \psi
+
\frac{1}{\sqrt{-g} }
\partial_{z}
\left[
\sqrt{- g}
\right]
g^{\mu\nu} \partial_{\mu} \partial_{\nu} \psi
\nn \\
-
\frac{1}{\sqrt{- g}}
\partial_{\mu}
\partial_{z}(\sqrt{-g} g^{\mu\nu}) \partial_{\nu} \psi
-
\frac{1}{\sqrt{- g}}
\partial_{z}(\sqrt{-g} g^{\mu\nu}) 
\partial_{\mu} \partial_{\nu} \psi
\nn \\
=
\partial_{z}
L\psi
+
\frac{1}{\sqrt{- g}}
\left[
\frac{\partial_{z}\sqrt{- g}}{\sqrt{-g}}
\partial_{\mu} (\sqrt{-g} g^{\mu\nu})
-
\partial_{\mu}
\partial_{z}(\sqrt{-g} g^{\mu\nu}) 
\right]
\partial_{\nu} \psi
-
\partial_{z} g^{\mu\nu}
\partial_{\mu} \partial_{\nu} \psi
\nn \\
=
\partial_{z}
L\psi
-
\frac{1}{\sqrt{- g}}
\left[
(\partial_{z}\sqrt{- g})
\Gamma^{\nu}_{\lambda \rho}
g^{\lambda \rho}
-
\partial_{z}(\sqrt{- g}\Gamma^{\nu}_{\lambda \rho}
g^{\lambda \rho}) 
\right]
\partial_{\nu} \psi
-
\partial_{z} g^{\mu\nu}
\partial_{\mu} \partial_{\nu} \psi
\nn \\
=
\partial_{z}
L\psi
-
\partial_{z} g^{\mu\nu}
\partial_{\mu} \partial_{\nu} \psi
+
\partial_{z}(\Gamma^{\lambda}_{\mu\nu}
g^{\mu\nu}) 
\partial_{\lambda} \psi
\end{gather}
%+++++++++++++++++++++++++++++++++++++++++++++++++++++++++++++++++++++++++++++++
Note that $\partial_{z}(\Gamma^{\lambda}_{\mu\nu} g^{\mu\nu})  =0$ as we showed earlier, and also
%+++++++++++++++++++++++++++++++++++++++++++++++++++++++++++++++++++++++++++++++
\begin{gather}
\partial_{z}g^{\mu\nu}
=
\left[
\begin{array}{cccc}
0 & 0 & 0 & 0
\\
0 & 0 & 0 & 0
\\
0 & 0 & 0 & - \frac{\bar{z} (1 + z \bar{z})}{r^{2}}
\\
0 & 0 & - \frac{\bar{z} (1 + z \bar{z})}{r^{2}} & 0
\end{array}
\right]
\end{gather}
%+++++++++++++++++++++++++++++++++++++++++++++++++++++++++++++++++++++++++++++++

%+++++++++++++++++++++++++++++++++++++++++++++++++++++++++++++++++++++++++++++++
\begin{gather}
L
\partial_{z}
\psi
=
\partial_{z}
L\psi
+
\frac{2 \bar{z} (1 + z \bar{z})}{r^{2}}
\partial_{\bar{z}} \partial_{z} \psi
\quad
\Rightarrow
\quad
\left(
L
-
\frac{2 \bar{z} (1 + z \bar{z})}{r^{2}}
\partial_{\bar{z}} 
\right)
\partial_{z} \psi
=
\partial_{z}
L\psi
\end{gather}
%+++++++++++++++++++++++++++++++++++++++++++++++++++++++++++++++++++++++++++++++
Looking at $L \partial_{\bar{z}}\psi$:
%+++++++++++++++++++++++++++++++++++++++++++++++++++++++++++++++++++++++++++++++
\begin{gather}
L \partial_{\bar{z}}\psi
=
\partial_{\bar{z}}
L\psi
-
\partial_{\bar{z}} g^{\mu\nu}
\partial_{\mu} \partial_{\nu} \psi
+
\partial_{\bar{z}}(\Gamma^{\lambda}_{\mu\nu}
g^{\mu\nu}) 
\partial_{\lambda} \psi
\end{gather}
%+++++++++++++++++++++++++++++++++++++++++++++++++++++++++++++++++++++++++++++++
Similarly,
%+++++++++++++++++++++++++++++++++++++++++++++++++++++++++++++++++++++++++++++++
\begin{gather}
L
\partial_{\bar{z}}
\psi
=
\partial_{\bar{z}}
L\psi
+
\frac{2 z (1 + z \bar{z})}{r^{2}}
\partial_{z} \partial_{\bar{z}} \psi
\quad
\Rightarrow
\quad
\left(
L
-
\frac{2 z (1 + z \bar{z})}{r^{2}}
\partial_{z} 
\right)
\partial_{\bar{z}} \psi
=
\partial_{\bar{z}}
L\psi
\end{gather}
%+++++++++++++++++++++++++++++++++++++++++++++++++++++++++++++++++++++++++++++++
The equations of motion are
%+++++++++++++++++++++++++++++++++++++++++++++++++++++++++++++++++++++++++++++++
\begin{gather}
L A_{r}
=
\frac{2}{r^{2}}
\partial_{r}
\left[
r
(1-
\frac{2 M}{r})
A_{r}
\right]
\label{SimpEqu1}
\\
\left(L 
-
\frac{2 \bar{z} (1 + z \bar{z})}{r^{2}}
\partial_{\bar{z}}
\right)
A_{z}
=
-
\frac{2}{r}
(1 - \frac{2 M}{r})
\partial_{r}
A_{z}
+
\frac{2}{r}
(1 - \frac{3 M}{r})
\partial_{z}
A_{r}
\label{SimpEqu2}
\\
\left(
L 
-
\frac{2 z (1 + z \bar{z})}{r^{2}}
\partial_{z}
\right)
A_{\bar{z}}
=
-
\frac{2}{r}
(1 - \frac{2 M}{r})
\partial_{r}
A_{\bar{z}}
+
\frac{2}{r}
(1 - \frac{3 M}{r})
\partial_{\bar{z}}
A_{r}
\label{SimpEqu3}
\end{gather}
%+++++++++++++++++++++++++++++++++++++++++++++++++++++++++++++++++++++++++++++++
and the constraint is
%+++++++++++++++++++++++++++++++++++++++++++++++++++++++++++++++++++++++++++++++
\begin{gather}
(1 - \frac{2 M}{r})
\partial_{r}
(r^{2}
A_{r})
+
\frac{(1 + z \bar{z})^{2}}{2}
(\partial_{z} A_{\bar{z}} + \partial_{\bar{z}} A_{z})
=0
\label{ConstraintEqu}
\end{gather}






%+++++++++++++++++++++++++++++++++++++++++++++++++++++++++++++++++++++++++++++++
%+++++++++++++++++++++++++++++++++++++++++++++++++++++++++++++++++++++++++++++++
%+++++++++++++++++++++++++++++++++++++++++++++++++++++++++++++++++++++++++++++++
%+++++++++++++++++++++++++++++++++++++++++++++++++++++++++++++++++++++++++++++++
\subsection{From Vector to Scalar Equations}
%+++++++++++++++++++++++++++++++++++++++++++++++++++++++++++++++++++++++++++++++
%+++++++++++++++++++++++++++++++++++++++++++++++++++++++++++++++++++++++++++++++
%+++++++++++++++++++++++++++++++++++++++++++++++++++++++++++++++++++++++++++++++

%+++++++++++++++++++++++++++++++++++++++++++++++++++++++++++++++++++++++++++++++
\subsubsection{First Solution}
%+++++++++++++++++++++++++++++++++++++++++++++++++++++++++++++++++++++++++++++++
%+++++++++++++++++++++++++++++++++++++++++++++++++++++++++++++++++++++++++++++++
%+++++++++++++++++++++++++++++++++++++++++++++++++++++++++++++++++++++++++++++++
%+++++++++++++++++++++++++++++++++++++++++++++++++++++++++++++++++++++++++++++++
%+++++++++++++++++++++++++++++++++++++++++++++++++++++++++++++++++++++++++++++++
Suppose $A_{z} = \partial_{z} \phi_{1}$, where $\phi_{1}$ is a scalar function, then equation (\ref{SimpEqu2}) becomes
%+++++++++++++++++++++++++++++++++++++++++++++++++++++++++++++++++++++++++++++++
\begin{gather}
\partial_{z}
L
\phi_{1}
=
-
\frac{2}{r}
(1 - \frac{2 M}{r})
\partial_{r}
\partial_{z}
\phi_{1}
+
\frac{2}{r}
(1 - \frac{3 M}{r})
\partial_{z}
A_{r}
\nn \\
\Rightarrow
\quad
\partial_{z}
\left(
L
\phi_{1}
+
\frac{2}{r}
(1 - \frac{2 M}{r})
\partial_{r}
\phi_{1}
-
\frac{2}{r}
(1 - \frac{3 M}{r})
A_{r}\right)
=0
\label{Phi1Equ}
\end{gather}
%+++++++++++++++++++++++++++++++++++++++++++++++++++++++++++++++++++++++++++++++
and if $A_{\bar{z}} = \partial_{\bar{z}} \phi_{2}$ then equation (\ref{SimpEqu3}) becomes
%+++++++++++++++++++++++++++++++++++++++++++++++++++++++++++++++++++++++++++++++
\begin{gather}
\partial_{\bar{z}}
L
\phi_{2}
=
-
\frac{2}{r}
(1 - \frac{2 M}{r})
\partial_{r}
\partial_{\bar{z}}
\phi_{2}
+
\frac{2}{r}
(1 - \frac{3 M}{r})
\partial_{\bar{z}}
A_{r}
\nn \\
\Rightarrow
\quad
\partial_{\bar{z}}
\left(
L
\phi_{2}
+
\frac{2}{r}
(1 - \frac{2 M}{r})
\partial_{r}
\phi_{2}
-
\frac{2}{r}
(1 - \frac{3 M}{r})
A_{r}\right)
=0
\label{Phi2Equ}
\end{gather}
%+++++++++++++++++++++++++++++++++++++++++++++++++++++++++++++++++++++++++++++++
Integrating and subtracting these two equations yields
%+++++++++++++++++++++++++++++++++++++++++++++++++++++++++++++++++++++++++++++++
\begin{gather}
L(\phi_{1} - \phi_{2})
+
\frac{2}{r}
(1 - \frac{2 M}{r})
\partial_{r}
(\phi_{1} -\phi_{2})
=
F(t,r,\bar{z})
-
G(t,r,z)
\end{gather}
%+++++++++++++++++++++++++++++++++++++++++++++++++++++++++++++++++++++++++++++++
where $F$ and $G$ are arbitrary functions that result from integration. The constraint equation (\ref{ConstraintEqu}) becomes  
%+++++++++++++++++++++++++++++++++++++++++++++++++++++++++++++++++++++++++++++++
\begin{gather}
(1 - \frac{2 M}{r})
\partial_{r}
(r^{2}
A_{r})
+
\frac{(1 + z \bar{z})^{2}}{2}
\partial_{z} \partial_{\bar{z}}( \phi_{1} + \phi_{2})
=0
\label{subtractEqu}
\end{gather}
%+++++++++++++++++++++++++++++++++++++++++++++++++++++++++++++++++++++++++++++++
What if $\phi_{1} = - \phi_{2} = \Phi$? Then the constraint becomes
%+++++++++++++++++++++++++++++++++++++++++++++++++++++++++++++++++++++++++++++++
\begin{gather}
(1 - \frac{2 M}{r})
\partial_{r}
(r^{2}
A_{r})
=0
\quad
\Rightarrow
\quad
r^{2} A_{r}
=
H(t,z,\bar{z})
\end{gather}
%+++++++++++++++++++++++++++++++++++++++++++++++++++++++++++++++++++++++++++++++
The $A_{r}$ equation of motion becomes
%+++++++++++++++++++++++++++++++++++++++++++++++++++++++++++++++++++++++++++++++
\begin{gather}
\frac{1}{1 - \frac{2 M}{r}}
\partial_{t}^{2}H
=
\frac{(1 + z \bar{z})^{2}}{r^{2}}
\partial_{z}
\partial_{\bar{z}}
H
\end{gather}
%+++++++++++++++++++++++++++++++++++++++++++++++++++++++++++++++++++++++++++++++
If $H$ is nonzero
%+++++++++++++++++++++++++++++++++++++++++++++++++++++++++++++++++++++++++++++++
\begin{gather}
\frac{\partial_{t}^{2}H}{\partial_{z}
\partial_{\bar{z}} H}
=
\frac{(1 + z \bar{z})^{2}}{r^{2} (\frac{1}{1 - \frac{2 M}{r}})}
\end{gather}
%+++++++++++++++++++++++++++++++++++++++++++++++++++++++++++++++++++++++++++++++
The left hand side is independent of $r$, thus we arrive at a contradiction, thus $H =0$. and equation (\ref{subtractEqu})
%+++++++++++++++++++++++++++++++++++++++++++++++++++++++++++++++++++++++++++++++
\begin{gather}
L\Phi
+
\frac{2}{r}
(1 - \frac{2 M}{r})
\partial_{r}
\Phi
=
\frac{F(t,r,\bar{z}) - G(t,r,z)}{2}
\end{gather}
%+++++++++++++++++++++++++++++++++++++++++++++++++++++++++++++++++++++++++++++++
Let's just say for now that the right hand side is zero then 
%+++++++++++++++++++++++++++++++++++++++++++++++++++++++++++++++++++++++++++++++
\begin{gather}
L\Phi
=
-\frac{2}{r}
(1 - \frac{2 M}{r})
\partial_{r}
\Phi
\end{gather}
%+++++++++++++++++++++++++++++++++++++++++++++++++++++++++++++++++++++++++++++++
This can be expressed as
%+++++++++++++++++++++++++++++++++++++++++++++++++++++++++++++++++++++++++++++++
\begin{gather}
g^{tt}
\partial_{t}^{2}\Phi
+
\frac{1}{r^{2}}
\partial_{r}
(r^{2} g^{rr} \partial_{r}\Phi)
+
(1 + z \bar{z})^{2}
\partial_{z}
((1 + z \bar{z})^{-2} g^{z \bar{z}} \partial_{\bar{z}} \Phi)
\nn \\
+
(1 + z \bar{z})^{2}
\partial_{\bar{z}}
((1 + z \bar{z})^{-2} g^{\bar{z} z} \partial_{z} \Phi)
=
-\frac{2}{r}
(1 - \frac{2 M}{r})
\partial_{r}
\Phi
\nn \\
\Rightarrow
\quad
(1 - \frac{2 M}{r})^{-1}
\partial_{t}^{2}\Phi
-
\frac{1}{r^{2}}
\partial_{r}
\left(
r^{2} (1 - \frac{2 M}{r}) \partial_{r}\Phi
\right)
-
\frac{(1 + z \bar{z})^{2}}{r^{2}}
\partial_{z}\partial_{\bar{z}} \Phi
=
-\frac{2}{r}
(1 - \frac{2 M}{r})
\partial_{r}
\Phi
\nn \\
\Rightarrow
\quad
(1 - \frac{2 M}{r})^{-1}
\partial_{t}^{2}\Phi
-
\frac{1}{r^{2}}
\partial_{r}
\left(
r^{2} (1 - \frac{2 M}{r}) \partial_{r}\Phi
\right)
\nn \\
-
\frac{1}{r^{2}}
\left(
\frac{1}{\sin\theta}
\frac{\partial}{\partial \theta}
\left(
\sin\theta
\frac{\partial \Phi}{\partial \theta}
\right)
+
\frac{1}{\sin^{2}\theta}
\frac{\partial^{2}\Phi}{\partial \varphi^{2}}
\right)
=
-\frac{2}{r}
(1 - \frac{2 M}{r})
\partial_{r}
\Phi
\end{gather}
%+++++++++++++++++++++++++++++++++++++++++++++++++++++++++++++++++++++++++++++++
Where we have used the fact that:
%+++++++++++++++++++++++++++++++++++++++++++++++++++++++++++++++++++++++++++++++
\begin{gather}
(1 + z \bar{z})^{2}
\partial_{z}\partial_{\bar{z}} \Phi
=
\frac{1}{\sin\theta}
\frac{\partial}{\partial \theta}
\left(
\sin\theta
\frac{\partial \Phi}{\partial \theta}
\right)
+
\frac{1}{\sin^{2}\theta}
\frac{\partial^{2}\Phi}{\partial \varphi^{2}}
\end{gather}
%+++++++++++++++++++++++++++++++++++++++++++++++++++++++++++++++++++++++++++++++
The proof is in the appendix. Let's suppose that $\Phi(t,r,\theta,\varphi) = f(r) Y_{l}^{m}(\theta,\varphi) e^{i \omega t}$. Thus
%+++++++++++++++++++++++++++++++++++++++++++++++++++++++++++++++++++++++++++++++
\begin{gather}
-(1 - \frac{2 M}{r})^{-1}
\omega^{2}f(r)
-
\frac{1}{r^{2}}
\partial_{r}
\left(
r^{2} (1 - \frac{2 M}{r}) \partial_{r}f(r)
\right)
+
\frac{l (l +1)}{r^{2}}f(r)
=
-\frac{2}{r}
(1 - \frac{2 M}{r})
\partial_{r}
f(r)
\nn \\
\Rightarrow
\quad
-(1 - \frac{2 M}{r})^{-1}
\omega^{2}f(r)
-
\frac{1}{r^{2}}
\partial_{r}
\left(
r^{2} (1 - \frac{2 M}{r}) \partial_{r}f(r)
\right)
+
\frac{2}{r}
(1 - \frac{2 M}{r})
\partial_{r}
f(r)
\nn \\
=
-
\frac{l (l +1)}{r^{2}}f(r)
\quad
\Rightarrow
\quad
-(1 - \frac{2 M}{r})^{-1}
\omega^{2}f(r)
-
\partial_{r}
\left(
(1 - \frac{2 M}{r}) \partial_{r}f(r)
\right)
\nn \\
-
\frac{2}{r}
(1 - \frac{2 M}{r}) \partial_{r}f(r)
+
\frac{2}{r}
(1 - \frac{2 M}{r})
\partial_{r}
f(r)
=
-
\frac{l (l +1)}{r^{2}}f(r)
\nn \\
\Rightarrow
\quad
(1 - \frac{2 M}{r})^{-1}
\omega^{2}f(r)
+
\partial_{r}
\left(
(1 - \frac{2 M}{r}) \partial_{r}f(r)
\right)
=
\frac{l (l +1)}{r^{2}}f(r)
\end{gather}
%+++++++++++++++++++++++++++++++++++++++++++++++++++++++++++++++++++++++++++++++
Define the tortoise coordinate as
%+++++++++++++++++++++++++++++++++++++++++++++++++++++++++++++++++++++++++++++++
\begin{gather}
r^{\ast}
=
r+
2 M\log\left|\frac{r}{2 M} -1\right|
\end{gather}
%+++++++++++++++++++++++++++++++++++++++++++++++++++++++++++++++++++++++++++++++
So
%+++++++++++++++++++++++++++++++++++++++++++++++++++++++++++++++++++++++++++++++
\begin{gather}
\frac{d r^{\ast}}{d r}
=
(1 - \frac{2 M}{r})^{-1}
\end{gather}
%+++++++++++++++++++++++++++++++++++++++++++++++++++++++++++++++++++++++++++++++
and therefore
%+++++++++++++++++++++++++++++++++++++++++++++++++++++++++++++++++++++++++++++++
\begin{gather}
(1 - \frac{2 M}{r})^{-1}
\omega^{2}f(r)
+
(1 - \frac{2 M}{r})^{-1}
\partial_{r^{\ast}}
\left(
\partial_{r^{\ast}}f(r)
\right)
=
\frac{l (l +1)}{r^{2}}f(r)
\end{gather}
%+++++++++++++++++++++++++++++++++++++++++++++++++++++++++++++++++++++++++++++++
Cleaning things up a bit
%+++++++++++++++++++++++++++++++++++++++++++++++++++++++++++++++++++++++++++++++
\begin{gather}
\left(
\partial_{r^{\ast}}^{2}
+
\omega^{2}
-
\frac{l (l +1)}{r^{2}}
(1 - \frac{2 M}{r})
\right)
f(r^{\ast})
=0
\end{gather}
%+++++++++++++++++++++++++++++++++++++++++++++++++++++++++++++++++++++++++++++++
the solution to this will give us the function for the gauge field which can be found from $\Phi$ by using
%+++++++++++++++++++++++++++++++++++++++++++++++++++++++++++++++++++++++++++++++
\begin{gather}
A_{\mu}
=
\left\{
0, 0, \partial_{z} \Phi , - \partial_{\bar{z}} \Phi
\right\}
\end{gather}
%+++++++++++++++++++++++++++++++++++++++++++++++++++++++++++++++++++++++++++++++




%+++++++++++++++++++++++++++++++++++++++++++++++++++++++++++++++++++++++++++++++
%+++++++++++++++++++++++++++++++++++++++++++++++++++++++++++++++++++++++++++++++
%+++++++++++++++++++++++++++++++++++++++++++++++++++++++++++++++++++++++++++++++
\subsubsection{A Second Solution}
%+++++++++++++++++++++++++++++++++++++++++++++++++++++++++++++++++++++++++++++++
%+++++++++++++++++++++++++++++++++++++++++++++++++++++++++++++++++++++++++++++++
%+++++++++++++++++++++++++++++++++++++++++++++++++++++++++++++++++++++++++++++++
Now suppose $\phi_{1} = \phi_{2} = \Psi$. Then (\ref{subtractEqu}) now becomes
%+++++++++++++++++++++++++++++++++++++++++++++++++++++++++++++++++++++++++++++++
\begin{gather}
(1 - \frac{2 M}{r})
\partial_{r}(r^{2} A_{r})
+
(1 + z \bar{z})^{2}
\partial_{z}\partial_{\bar{z}}\Psi
=0
\end{gather}
%+++++++++++++++++++++++++++++++++++++++++++++++++++++++++++++++++++++++++++++++
and adding equations (\ref{Phi1Equ}) and (\ref{Phi2Equ}) we get
%+++++++++++++++++++++++++++++++++++++++++++++++++++++++++++++++++++++++++++++++
\begin{gather}
L \Psi
+
\frac{2}{r}
(1 - \frac{2 M}{r})
\partial_{r}\Psi
-
\frac{2}{r}
\left(
1
-
\frac{3 M}{r}
\right)
A_{r}
=0
\nn \\
\Rightarrow
\quad
(1 - \frac{2 M}{r})^{-1} \partial^{2}_{t}
\Psi
-
\frac{1}{r^{2}}
\partial_{r}
\left(
r^{2}
(1 - \frac{2 M}{r})
\partial_{r} \Psi
\right)
-
\frac{(1 + z \bar{z})^{2}}{r^{2}}
\partial_{z}
\partial_{\bar{z}}
\Psi
\nn \\
+
\frac{2}{r}
(1 - \frac{2 M}{r})
\partial_{r}\Psi
-
\frac{2}{r}
\left(
1
-
\frac{3 M}{r}
\right)
A_{r}
=0
\nn \\
\Rightarrow
\quad
(1 - \frac{2 M}{r})^{-1} \partial^{2}_{t}
\Psi
-
\partial_{r}
\left(
(1 - \frac{2 M}{r})
\partial_{r} \Psi
\right)
-
\frac{(1 + z \bar{z})^{2}}{r^{2}}
\partial_{z}
\partial_{\bar{z}}
\Psi
-
\frac{2}{r}
\left(
1
-
\frac{3 M}{r}
\right)
A_{r}
=0
\end{gather}
%+++++++++++++++++++++++++++++++++++++++++++++++++++++++++++++++++++++++++++++++
Let $\Psi(t,r,\theta,\varphi) = (1 - \frac{2 M}{r}) \partial_{r}\tilde{\Phi}(r) Y_{l}^{m}(\theta,\varphi) e^{i \omega t}$
%+++++++++++++++++++++++++++++++++++++++++++++++++++++++++++++++++++++++++++++++
\begin{gather}
-
\omega^{2}
\partial_{r}
\tilde{\Phi}(r)
-
\partial_{r}
\left(
(1 - \frac{2 M}{r})
\partial_{r} \left[ (1 - \frac{2 M}{r}) \partial_{r}\tilde{\Phi}(r)\right]
\right)
\nn \\
+
(1 - \frac{2 M}{r})
\frac{l (l +1)}{r^{2}}
\partial_{r}
\tilde{\Phi}(r)
-
\frac{2}{r}
\left(
1
-
\frac{3 M}{r}
\right)
\frac{A_{r} e^{- i \omega t}}{Y_{l}^{m}}
=0
\nn \\
\Rightarrow
\quad
-
\partial_{r}
\left[
\left(
(1 - \frac{2 M}{r})
\partial_{r} \left[ (1 - \frac{2 M}{r}) \partial_{r}\tilde{\Phi}\right]
\right)
+
\omega^{2}
\tilde{\Phi}
\right]
\nn \\
+
(1 - \frac{2 M}{r})
\frac{l (l +1)}{r^{2}}
\partial_{r}
\tilde{\Phi}
-
\frac{2}{r}
\left(
1
-
\frac{3 M}{r}
\right)
\frac{A_{r} e^{- i \omega t}}{Y_{l}^{m}}
=0
\nn \\
\Rightarrow
\quad
-
\partial_{r}
\left[
\partial_{r^{\ast}}^{2}
\tilde{\Phi}
+
\omega^{2}
\tilde{\Phi}
\right]
+
\partial_{r}
\left[
(1 - \frac{2 M}{r})
\frac{l (l +1)}{r^{2}}
\tilde{\Phi}
\right]
\nn \\
-
\partial_{r}
\left[
(1 - \frac{2 M}{r})
\frac{l (l +1)}{r^{2}}
\right]
\tilde{\Phi}
-
\frac{2}{r}
\left(
1
-
\frac{3 M}{r}
\right)
\frac{A_{r} e^{- i \omega t}}{Y_{l}^{m}}
=0
\nn \\
\Rightarrow
\quad
-
\partial_{r}
\left[
\partial^{2}_{r^{\ast}}
\tilde{\Phi}
+
\left(
\omega^{2}
-
(1 - \frac{2 M}{r})
\frac{l (l +1)}{r^{2}}
\right)
\tilde{\Phi}
\right]
\nn \\
+
\frac{2}{r}
\left(1 - \frac{3 M}{r}\right)
\left[
\frac{l (l + 1)}{r^{2}}
\tilde{\Phi}
-
\frac{A_{r} e^{- i \omega t}}{Y_{l}^{m}}
\right]
=0
\end{gather}
%+++++++++++++++++++++++++++++++++++++++++++++++++++++++++++++++++++++++++++++++
This Is simplified if $A_{r}(t,r,\theta,\varphi) = l (l + 1) \frac{\tilde{\Phi}(r)}{r^{2}} Y_{l}^{m}(\theta,\varphi) e^{i \omega t}$, in which case the equation of motion becomes
%+++++++++++++++++++++++++++++++++++++++++++++++++++++++++++++++++++++++++++++++
\begin{gather}
\partial_{r}
\left[
\partial^{2}_{r^{\ast}}
\tilde{\Phi}
+
\left(
\omega^{2}
-
(1 - \frac{2 M}{r})
\frac{l (l +1)}{r^{2}}
\right)
\tilde{\Phi}
\right]
=0
\nn \\
\Rightarrow
\quad
\left(
\partial^{2}_{r^{\ast}}
+
\omega^{2}
-
\frac{l (l +1)}{r^{2}}
(1 - \frac{2 M}{r})
\right)
\tilde{\Phi}
=0
\end{gather}
%+++++++++++++++++++++++++++++++++++++++++++++++++++++++++++++++++++++++++++++++
Let's call $\Phi(t,r,\theta,\varphi) = \tilde{\Phi}(r) Y_{l}^{m}(\theta,\varphi) e^{i \omega t}$. The second solution for the EM field is given by.
%+++++++++++++++++++++++++++++++++++++++++++++++++++++++++++++++++++++++++++++++
\begin{gather}
A_{\mu}
=
\left\{
0, \frac{l (l+1)}{r^{2}}
\Phi
,
(1 - \frac{2 M}{r})
\partial_{z}\partial_{r}\Phi
,
(1 - \frac{2 M}{r})
\partial_{\bar{z}}\partial_{r}\Phi
\right\}
\end{gather}
%+++++++++++++++++++++++++++++++++++++++++++++++++++++++++++++++++++++++++++++++





































%\bibliographystyle{h-elsevier}
%\bibliography{Neutrino-K-Essence}





\end{document}